\section{Spectrally Preconditioned Truncated-SVD HPAO Solver}

\subsection{Постановка}

Рассматривается непосредственно rank-constrained global step
multi-index HPAO.

При фиксированном текущем basis
\[
    P\in\mathbb R^{m\times d},
    \qquad
    PP^\top=I_m,
\]
и фиксированных локальных коэффициентах
\[
    g_j\in\mathbb R^m
\]
необходимо решить
\[
\boxed{
    \min_{\rank(B)\le r}
    F(B)
}
\]
для
\[
\boxed{
    F(B)
    =
    \sum_{j=1}^J
    N_j
    \left\|
        I_j-U_jB^\top g_j
    \right\|_2^2
    +
    \lambda
    \|B-P\|_F^2.
}
\]

Здесь
\[
    B\in\mathbb R^{m\times d},
    \qquad
    r<m,
\]
и предполагается
\[
    \lambda>0.
\]

Основная задача состоит именно в построении малорангового
приближения самой матрицы $B$:
\[
\boxed{
    B_r
    =
    A\Sigma V^\top,
    \qquad
    A^\top A=I_r,
    \quad
    V^\top V=I_r.
}
\]

Не вводится представление
\[
    B=P+\Delta.
\]

\subsection{Действительно различные семейства методов}

Для данной rank-constrained задачи следует различать как минимум
четыре математически различных семейства подходов.

\paragraph{I. Greedy conditional spectral recovery.}

Матрица строится последовательно:
\[
    B_k
    =
    B_{k-1}
    +
    \sigma_k a_kv_k^\top.
\]

На каждом шаге выбирается rank-one направление,
дающее максимальное уменьшение исходного квадратичного functional.

Для фиксированного rank-one направления optimal line-search
получается аналитически.

Глобальный поиск лучшего $(a,v)$ для общего набора $H_j$
остаётся нелинейной биквадратичной задачей.

\paragraph{II. Exact whitening + truncated SVD.}

При специальной структуре
\[
    H_j=\eta_jI_d
\]
вся rank-$r$ задача точно преобразуется в best rank-$r$
matrix approximation в норме Фробениуса.

В этом случае truncated SVD даёт глобальное решение исходной
rank-constrained задачи.

\paragraph{III. Block low-rank factor optimization.}

Используется
\[
    B=XY^\top,
\]
где
\[
    X\in\mathbb R^{m\times r},
    \qquad
    Y\in\mathbb R^{d\times r},
\]
и выполняется alternating least squares.

Этот подход не требует изотропии $H_j$ и одновременно оптимизирует
все $r$ компонент, но задача совместно по $(X,Y)$ невыпукла.

\paragraph{IV. Convex nuclear-norm relaxation.}

Ограничение
\[
    \rank(B)\le r
\]
можно заменить penalty
\[
    \tau\|B\|_*.
\]

Это приводит к другой оптимизационной задаче.
Она может использоваться как baseline, но не является точным
решением исходной rank-constrained постановки.

\subsection{Квадратичная операторная форма}

Определим
\[
    H_j
    =
    U_j^\top U_j
    \succeq0
\]
и
\[
    c_j
    =
    U_j^\top I_j.
\]

Раскрывая квадрат:
\[
\begin{aligned}
    \|I_j-U_jB^\top g_j\|_2^2
    ={}&
    \|I_j\|_2^2
    -
    2g_j^\top Bc_j
    \\
    &+
    g_j^\top BH_jB^\top g_j.
\end{aligned}
\]

Также
\[
\begin{aligned}
    \lambda\|B-P\|_F^2
    ={}&
    \lambda\|B\|_F^2
    -
    2\lambda\langle B,P\rangle_F
    +
    \lambda\|P\|_F^2.
\end{aligned}
\]

Следовательно,
\[
\boxed{
    F(B)
    =
    C_0
    +
    \langle B,\mathcal L(B)\rangle_F
    -
    2\langle C_B,B\rangle_F,
}
\]
где
\[
\boxed{
    \mathcal L(X)
    =
    \sum_{j=1}^J
    N_j
    (g_jg_j^\top)
    XH_j
    +
    \lambda X
}
\]
и
\[
\boxed{
    C_B
    =
    \sum_{j=1}^J
    N_jg_jc_j^\top
    +
    \lambda P.
}
\]

Константа равна
\[
    C_0
    =
    \sum_jN_j\|I_j\|_2^2
    +
    \lambda\|P\|_F^2.
\]

Так как
\[
    PP^\top=I_m,
\]
имеем
\[
    \|P\|_F^2=m.
\]

\subsection{Лемма 1: свойства оператора $\mathcal L$}

\textbf{Лемма.}
Оператор $\mathcal L$ самосопряжён и положительно определён.
Для любого
\[
    X\in\mathbb R^{m\times d}
\]
выполняется
\[
\boxed{
    \langle X,\mathcal L(X)\rangle_F
    =
    \sum_j
    N_j
    \|U_jX^\top g_j\|_2^2
    +
    \lambda\|X\|_F^2
    \ge
    \lambda\|X\|_F^2.
}
\]

\textbf{Доказательство.}

Для каждого $j$
\[
\begin{aligned}
    \left\langle
        X,
        (g_jg_j^\top)XH_j
    \right\rangle_F
    &=
    \operatorname{tr}
    \left(
        X^\top g_jg_j^\top XU_j^\top U_j
    \right)
    \\
    &=
    \|U_jX^\top g_j\|_2^2.
\end{aligned}
\]

Отсюда следует положительная определённость при $\lambda>0$.

Самосопряжённость следует из симметричности
\[
    g_jg_j^\top
    \quad\text{и}\quad
    H_j.
\]

\subsection{Полное unconstrained решение}

Без ограничения rank задача строго выпукла.

Её единственный minimizer $B_\star$ удовлетворяет
\[
\boxed{
    \mathcal L(B_\star)=C_B.
}
\]

Кроме того,
\[
\boxed{
    F(B)
    =
    F(B_\star)
    +
    \langle
        B-B_\star,
        \mathcal L(B-B_\star)
    \rangle_F.
}
\]

Следовательно, исходная rank-constrained задача эквивалентна
\[
\boxed{
    B_r^\star
    =
    \arg\min_{\rank(B)\le r}
    \|B-B_\star\|_{\mathcal L}^2,
}
\]
где
\[
    \|X\|_{\mathcal L}^2
    =
    \langle X,\mathcal L(X)\rangle_F.
\]

Таким образом, задача является generalized low-rank approximation
в Hessian-метрике исходного HPAO functional.

\subsection{Normal residual}

Для текущего приближения $B_t$ определим
\[
\boxed{
    Q_t
    =
    C_B-\mathcal L(B_t).
}
\]

Тогда
\[
\boxed{
    \nabla F(B_t)=-2Q_t.
}
\]

Для произвольной поправки $Z$
имеем точное равенство
\[
\boxed{
    F(B_t+Z)-F(B_t)
    =
    \langle Z,\mathcal L(Z)\rangle_F
    -
    2\langle Q_t,Z\rangle_F.
}
\]

\subsection{Теорема 1: лучший rank-one line-search}

Пусть
\[
    Z=av^\top,
    \qquad
    a\neq0,
    \quad
    v\neq0.
\]

Рассмотрим
\[
    B_t+\sigma Z.
\]

Тогда
\[
    F(B_t+\sigma Z)-F(B_t)
    =
    D(a,v)\sigma^2
    -
    2R_t(a,v)\sigma,
\]
где
\[
\boxed{
    R_t(a,v)
    =
    a^\top Q_tv
}
\]
и
\[
\boxed{
    D(a,v)
    =
    \sum_j
    N_j
    (a^\top g_j)^2
    v^\top H_jv
    +
    \lambda
    \|a\|_2^2\|v\|_2^2.
}
\]

Так как
\[
    D(a,v)>0,
\]
optimal scale равен
\[
\boxed{
    \sigma^\star(a,v)
    =
    \frac{
        a^\top Q_tv
    }{
        D(a,v)
    }.
}
\]

Точное уменьшение objective:
\[
\boxed{
    \operatorname{gain}_t(a,v)
    =
    \frac{
        (a^\top Q_tv)^2
    }{
        D(a,v)
    }.
}
\]

Следовательно, лучший следующий rank-one шаг задаётся
\[
\boxed{
    \Gamma_1(Q_t)
    =
    \sup_{a,v\neq0}
    \frac{
        (a^\top Q_tv)^2
    }{
        \sum_j
        N_j
        (a^\top g_j)^2
        v^\top H_jv
        +
        \lambda
        \|a\|^2\|v\|^2
    }.
}
\]

\subsection{Conditional spectral norm}

Определим
\[
\boxed{
    \|Q\|_{\mathrm{csp},\mathcal L}
    =
    \sup_{\rank(Z)=1}
    \frac{
        |\langle Q,Z\rangle_F|
    }{
        \sqrt{
            \langle Z,\mathcal L(Z)\rangle_F
        }
    }.
}
\]

Тогда
\[
\boxed{
    \Gamma_1(Q)
    =
    \|Q\|_{\mathrm{csp},\mathcal L}^2.
}
\]

Это даёт точный смысл spectral criterion:
\[
\boxed{
    \text{conditional spectral norm}^2
    =
    \text{максимально возможный gain одного rank-one шага}.
}
\]

\subsection{Conditional optimization по $a$}

При фиксированном единичном $v$ определим
\[
\boxed{
    G(v)
    =
    \sum_j
    N_j
    (v^\top H_jv)
    g_jg_j^\top
    +
    \lambda I_m.
}
\]

Тогда
\[
    D(a,v)
    =
    a^\top G(v)a
\]
и
\[
    R_t(a,v)
    =
    a^\top Q_tv.
\]

Следовательно,
\[
\boxed{
    \max_{a\neq0}
    \operatorname{gain}_t(a,v)
    =
    (Q_tv)^\top
    G(v)^{-1}
    (Q_tv),
}
\]
и optimal direction:
\[
\boxed{
    a
    \propto
    G(v)^{-1}Q_tv.
}
\]

\subsection{Conditional optimization по $v$}

При фиксированном единичном $a$ определим
\[
\boxed{
    H(a)
    =
    \sum_j
    N_j
    (a^\top g_j)^2H_j
    +
    \lambda I_d.
}
\]

Тогда
\[
\boxed{
    v
    \propto
    H(a)^{-1}Q_t^\top a.
}
\]

Таким образом, exact conditional spectral direction может
искаться alternating iterations
\[
\boxed{
    v
    \rightarrow
    a
    \rightarrow
    v
    \rightarrow
    \cdots.
}
\]

Каждый точный conditional update не уменьшает
$\operatorname{gain}_t(a,v)$.

Однако для произвольных $H_j$ совместная задача остаётся невыпуклой.
Глобальная оптимальность alternating procedure не гарантируется.

\subsection{Обусловленность conditional solves}

При
\[
    \|a\|=\|v\|=1
\]
имеем
\[
    G(v)\succeq\lambda I_m
\]
и
\[
    H(a)\succeq\lambda I_d.
\]

Следовательно,
\[
\boxed{
    \kappa(G(v))
    \le
    1+
    \frac{
        \sum_j
        N_j
        \|g_j\|_2^2
        \|H_j\|_2
    }{
        \lambda
    }
}
\]
и аналогично
\[
\boxed{
    \kappa(H(a))
    \le
    1+
    \frac{
        \sum_j
        N_j
        \|g_j\|_2^2
        \|H_j\|_2
    }{
        \lambda
    }.
}
\]

\subsection{Точная изотропная модель}

Теперь предположим
\[
\boxed{
    H_j=\eta_jI_d,
    \qquad
    \eta_j\ge0.
}
\]

Определим
\[
\boxed{
    G
    =
    \sum_j
    N_j\eta_jg_jg_j^\top
    +
    \lambda I_m.
}
\]

Тогда
\[
\boxed{
    \mathcal L(X)=GX.
}
\]

Линейная часть при этом остаётся
\[
\boxed{
    C_B
    =
    \sum_j
    N_jg_jc_j^\top
    +
    \lambda P.
}
\]

Это принципиальное отличие pure-$B$ постановки:
prior $P$ входит непосредственно в spectral target.

\subsection{Теорема 2: точное truncated-SVD решение}

Пусть
\[
    H_j=\eta_jI_d.
\]

Определим
\[
\boxed{
    M
    =
    G^{-1/2}C_B.
}
\]

Пусть
\[
    M
    =
    U
    \diag(s_1,\ldots,s_m)
    V^\top,
\]
где
\[
    s_1\ge s_2\ge\cdots\ge0.
\]

Тогда глобальный minimizer
\[
    \min_{\rank(B)\le r}F(B)
\]
равен
\[
\boxed{
    B_r^\star
    =
    G^{-1/2}
    [G^{-1/2}C_B]_r.
}
\]

\textbf{Доказательство.}

При изотропии
\[
\begin{aligned}
    F(B)
    &=
    C_0
    +
    \operatorname{tr}(B^\top GB)
    -
    2\operatorname{tr}(C_B^\top B).
\end{aligned}
\]

Введём
\[
    Z=G^{1/2}B.
\]

Так как $G$ положительно определена,
\[
    \rank(Z)=\rank(B).
\]

Тогда
\[
\begin{aligned}
    F(B)
    &=
    C_0
    +
    \|Z\|_F^2
    -
    2
    \langle
        G^{-1/2}C_B,
        Z
    \rangle_F
    \\
    &=
    C_0
    -
    \|M\|_F^2
    +
    \|Z-M\|_F^2.
\end{aligned}
\]

По теореме Eckart--Young
глобально оптимальное rank-$r$ решение есть
\[
    Z_r^\star=[M]_r.
\]

Следовательно,
\[
\boxed{
    B_r^\star
    =
    G^{-1/2}[M]_r.
}
\]

\subsection{Точное значение objective gap}

В тех же условиях
\[
\boxed{
    F(B_r^\star)
    =
    C_0
    -
    \sum_{k=1}^r s_k^2.
}
\]

Unconstrained solution:
\[
\boxed{
    B_\star
    =
    G^{-1}C_B.
}
\]

Для него
\[
\boxed{
    F(B_\star)
    =
    C_0
    -
    \sum_k s_k^2.
}
\]

Поэтому
\[
\boxed{
    F(B_r^\star)
    -
    F(B_\star)
    =
    \sum_{k>r}s_k^2.
}
\]

Таким образом, spectral tail имеет точный смысл ошибки
rank-$r$ approximation по исходному HPAO objective.

\subsection{Точный выбор rank}

Определим captured objective energy:
\[
\boxed{
    E_r
    =
    \frac{
        \sum_{k=1}^r s_k^2
    }{
        \sum_k s_k^2
    }.
}
\]

Тогда выбор
\[
\boxed{
    r_\varepsilon
    =
    \min
    \left\{
        r:
        \frac{
            \sum_{k>r}s_k^2
        }{
            \sum_ks_k^2
        }
        \le
        \varepsilon_{\mathrm{energy}}
    \right\}
}
\]
точно гарантирует
\[
    E_{r_\varepsilon}
    \ge
    1-\varepsilon_{\mathrm{energy}}.
\]

\subsection{Теорема 3: следующий singular component}

Для
\[
    B_r^\star
    =
    G^{-1/2}[M]_r
\]
normal residual равен
\[
    Q_r
    =
    C_B-GB_r^\star.
\]

Следовательно,
\[
\boxed{
    G^{-1/2}Q_r
    =
    M-[M]_r.
}
\]

Поэтому
\[
\boxed{
    \|Q_r\|_{\mathrm{csp},\mathcal L}
    =
    \|G^{-1/2}Q_r\|_2
    =
    s_{r+1}.
}
\]

Точный лучший следующий rank-one gain:
\[
\boxed{
    \Gamma_1(Q_r)=s_{r+1}^2.
}
\]

Более того, лучший rank-$q$ expansion имеет вид
\[
\boxed{
    \delta_q^\star
    =
    G^{-1/2}
    [M-[M]_r]_q
}
\]
и даёт улучшение
\[
\boxed{
    F(B_r^\star)
    -
    F(B_r^\star+\delta_q^\star)
    =
    \sum_{k=r+1}^{r+q}s_k^2.
}
\]

Таким образом, в точной изотропной модели нет необходимости
жадно угадывать следующую singular component:
весь оптимальный блок top-$q$ находится одним truncated SVD.

\subsection{Изотропная инициализация для общего случая}

В общем случае
\[
    H_j\neq\eta_jI.
\]

Выберем
\[
\boxed{
    \eta_j
    =
    \frac{\operatorname{tr}(H_j)}{d}
    =
    \frac{\|U_j\|_F^2}{d}.
}
\]

Тогда
\[
    \eta_jI
\]
является лучшей scalar approximation $H_j$
в норме Фробениуса:
\[
\boxed{
    \eta_j
    =
    \arg\min_\eta
    \|H_j-\eta I\|_F^2.
}
\]

Построим
\[
\boxed{
    G
    =
    \sum_j
    N_j\eta_jg_jg_j^\top
    +
    \lambda I_m
}
\]
и
\[
\boxed{
    C_B
    =
    \sum_j
    N_jg_jc_j^\top
    +
    \lambda P.
}
\]

Тогда естественная spectral initialization:
\[
\boxed{
    B_r^{(0)}
    =
    G^{-1/2}
    [G^{-1/2}C_B]_r.
}
\]

В общем случае это не является точным решением исходной задачи,
но может использоваться как:

\begin{enumerate}
    \item начальное приближение для greedy solver;
    \item начальное приближение для block ALS;
    \item deterministic fallback;
    \item источник начального rank;
    \item spectral preconditioner для residual-gradient.
\end{enumerate}

\subsection{Perturbation относительно изотропной модели}

Положим
\[
    E_j
    =
    H_j-\eta_jI
\]
и
\[
\boxed{
    \rho
    =
    \sum_j
    N_j
    \|g_j\|_2^2
    \|E_j\|_2.
}
\]

Пусть
\[
    \widetilde{\mathcal L}(X)=GX.
\]

Тогда
\[
\boxed{
    \|
        (\mathcal L-\widetilde{\mathcal L})(X)
    \|_F
    \le
    \rho\|X\|_F.
}
\]

Если
\[
    \rho<\lambda
\]
и
\[
    \theta=\rho/\lambda,
\]
то
\[
\boxed{
    (1-\theta)
    \langle X,\widetilde{\mathcal L}(X)\rangle_F
    \le
    \langle X,\mathcal L(X)\rangle_F
    \le
    (1+\theta)
    \langle X,\widetilde{\mathcal L}(X)\rangle_F.
}
\]

\subsection{Изотропная spectral proposal для общего residual}

Пусть текущая матрица равна $B_t$ и вычислен exact normal residual
\[
\boxed{
    Q_t
    =
    C_B-\mathcal L(B_t).
}
\]

Вместо обычного top singular vector $Q_t$
рассмотрим preconditioned matrix
\[
\boxed{
    W_t
    =
    G^{-1/2}Q_t.
}
\]

Её top singular pair решает isotropic surrogate задачи
лучшего rank-one шага.

Если
\[
    \rho<\lambda,
\]
то exact gain этого isotropic-optimal направления
составляет не менее
\[
\boxed{
    \frac{1-\theta}{1+\theta}
}
\]
доли от глобально лучшего exact rank-one gain.

Таким образом,
\[
\boxed{
    G^{-1/2}Q_t
}
\]
является curvature-aware заменой обычного residual-gradient
для инициализации следующей singular component.

\subsection{Block spectral proposal}

Вместо одной компоненты вычислить
\[
\boxed{
    [G^{-1/2}Q_t]_q
    =
    \widetilde A_q
    S_q
    V_q^\top.
}
\]

Использовать
\[
\boxed{
    Z_q^{(0)}
    =
    G^{-1/2}
    \widetilde A_q
    S_q
    V_q^\top
}
\]
как начальное rank-$q$ направление.

При точной изотропии это exact optimal rank-$q$ residual correction.

В общем случае это preconditioned spectral proposal,
который необходимо затем дооптимизировать по исходному functional.

\subsection{Block ALS непосредственно для $B$}

Вместо жадного восстановления по одной компоненте
рассмотреть факторизацию
\[
\boxed{
    B=XY^\top,
}
\]
где
\[
    X\in\mathbb R^{m\times r},
    \qquad
    Y\in\mathbb R^{d\times r}.
\]

Тогда
\[
    B^\top g_j
    =
    Y(X^\top g_j).
\]

Исходный functional принимает вид
\[
\boxed{
\begin{aligned}
    F(X,Y)
    ={}&
    \sum_j
    N_j
    \left\|
        I_j-U_jY(X^\top g_j)
    \right\|_2^2
    \\
    &+
    \lambda
    \|XY^\top-P\|_F^2.
\end{aligned}
}
\]

При фиксированном $X$ это convex quadratic problem по $Y$.

При фиксированном $Y$ это convex quadratic problem по $X$.

\subsection{Монотонность block ALS}

Если обе block subproblems решаются точно:
\[
    Y_{t+1}
    \in
    \arg\min_YF(X_t,Y),
\]
\[
    X_{t+1}
    \in
    \arg\min_XF(X,Y_{t+1}),
\]
то
\[
\boxed{
    F(X_{t+1},Y_{t+1})
    \le
    F(X_t,Y_t).
}
\]

Следовательно, objective sequence монотонна и ограничена снизу.

Это не доказывает глобальную сходимость к
\[
    B_r^\star,
\]
так как parameterized problem по $(X,Y)$ совместно невыпукла.

\subsection{Spectral initialization block ALS}

Предлагается инициализировать block ALS не случайно, а через
\[
\boxed{
    B_r^{(0)}
    =
    G^{-1/2}
    [G^{-1/2}C_B]_r.
}
\]

Для получения balanced factors выполнить SVD
\[
    B_r^{(0)}
    =
    A\Sigma V^\top
\]
и положить
\[
\boxed{
    X_0
    =
    A\Sigma^{1/2},
    \qquad
    Y_0
    =
    V\Sigma^{1/2}.
}
\]

Это предотвращает произвольный масштаб между факторами.

\subsection{Balancing после ALS}

После каждого полного ALS sweep использовать compact SVD:
\[
    X_tY_t^\top
    =
    A\Sigma V^\top
\]
и переопределять
\[
\boxed{
    X_t
    \leftarrow
    A\Sigma^{1/2},
    \qquad
    Y_t
    \leftarrow
    V\Sigma^{1/2}.
}
\]

Это не меняет $B$ и значение objective.

\subsection{Критерии выбора rank}

Следует различать два spectral criteria.

\paragraph{Energy criterion.}

Для singular values isotropic target
\[
    M=G^{-1/2}C_B
\]
использовать
\[
\boxed{
    \frac{
        \sum_{k>r}s_k^2
    }{
        \sum_ks_k^2
    }
    \le
    \varepsilon_{\mathrm{energy}}.
}
\]

При точной изотропии это exact relative objective-gap criterion.

При общей структуре $H_j$ это только surrogate.

\paragraph{Next-component criterion.}

Использовать
\[
\boxed{
    \frac{s_{r+1}}{s_1}
    \le
    \varepsilon_{\mathrm{sv}}.
}
\]

При точной изотропии
\[
\boxed{
    \frac{
        \text{best next rank-one gain}
    }{
        \text{best first rank-one gain}
    }
    =
    \left(
        \frac{s_{r+1}}{s_1}
    \right)^2.
}
\]

Однако малый $s_{r+1}$ сам по себе не контролирует сумму всего
оставшегося spectral tail.

Поэтому для выбора rank предпочтителен energy criterion,
а $\varepsilon_{\mathrm{sv}}$ следует использовать как тест
отсутствия одной сильной следующей компоненты.

\subsection{Spectral residual stopping}

Для текущего $B_t$ вычислить exact residual
\[
    Q_t=C_B-\mathcal L(B_t).
\]

Затем
\[
    W_t=G^{-1/2}Q_t.
\]

В точной изотропной модели
\[
\boxed{
    \|W_t\|_2^2
}
\]
равно максимально возможному улучшению одним rank-one шагом.

При perturbation condition
\[
    \rho<\lambda
\]
имеем upper certificate
\[
\boxed{
    \Gamma_1(Q_t)
    \le
    \frac{
        \|G^{-1/2}Q_t\|_2^2
    }{
        1-\rho/\lambda
    }.
}
\]

Следовательно, можно остановиться при
\[
\boxed{
    \frac{
        \|G^{-1/2}Q_t\|_2^2
    }{
        1-\rho/\lambda
    }
    \le
    \varepsilon_{\mathrm{gain}}.
}
\]

Тогда никакая дополнительная rank-one component не может уменьшить
исходный objective более чем на
$\varepsilon_{\mathrm{gain}}$.

\subsection{Предлагаемый pure truncated-SVD алгоритм}

Предлагается следующий алгоритм.

\paragraph{Шаг 1. Построение isotropic model.}

За один проход вычислить
\[
    c_j=U_j^\top I_j
\]
и
\[
    \eta_j
    =
    \frac{\|U_j\|_F^2}{d}.
\]

Сформировать
\[
    C_B
    =
    \sum_j
    N_jg_jc_j^\top
    +
    \lambda P
\]
и
\[
    G
    =
    \sum_j
    N_j\eta_jg_jg_j^\top
    +
    \lambda I_m.
\]

\paragraph{Шаг 2. Spectral initialization.}

Вычислить
\[
    M=G^{-1/2}C_B
\]
и truncated SVD
\[
    [M]_r.
\]

Положить
\[
\boxed{
    B_0
    =
    G^{-1/2}[M]_r.
}
\]

Rank $r$ можно выбрать автоматически по
$\varepsilon_{\mathrm{energy}}$.

\paragraph{Шаг 3. Exact block refinement.}

Факторизовать
\[
    B_0=X_0Y_0^\top
\]
и выполнить несколько block ALS sweeps,
оптимизируя полный исходный functional
\[
    F(B),
\]
а не isotropic surrogate.

\paragraph{Шаг 4. Exact normal residual.}

После refinement вычислить
\[
    Q_t
    =
    C_B-\mathcal L(B_t).
\]

\paragraph{Шаг 5. Spectral expansion.}

Если residual остаётся значимым,
вычислить
\[
    [G^{-1/2}Q_t]_q.
\]

Использовать соответствующий rank-$q$ block
для расширения текущего low-rank представления.

\paragraph{Шаг 6. Compression.}

После расширения выполнить compact SVD и сохранить rank budget.

\paragraph{Шаг 7. Exact refinement.}

Снова выполнить block ALS либо conditional
rank-one/rank-$q$ refinement по исходному functional.

\paragraph{Шаг 8. Stopping.}

Остановить solver по комбинации:

\[
    \text{objective gain},
\]
\[
    \varepsilon_{\mathrm{sv}},
\]
\[
    \varepsilon_{\mathrm{energy}},
\]
и, когда доступен perturbation certificate,
\[
    \|G^{-1/2}Q_t\|_2.
\]

\subsection{Вычислительная стоимость spectral initialization}

Для каждого $j$ необходимо вычислить
\[
    U_j^\top I_j
\]
и
\[
    \|U_j\|_F^2.
\]

При
\[
    U_j\in\mathbb R^{p\times d}
\]
это требует
\[
\boxed{
    O(Jpd)
}
\]
операций.

Accumulation
\[
    C_B
    =
    \sum_jN_jg_jc_j^\top+\lambda P
\]
требует
\[
    O(Jmd).
\]

Построение $G$:
\[
    O(Jm^2).
\]

Factorization $G$:
\[
    O(m^3).
\]

SVD матрицы
\[
    G^{-1/2}C_B\in\mathbb R^{m\times d}
\]
при $m\ll d$:
\[
    O(m^2d).
\]

Итого:
\[
\boxed{
    O(
        Jpd
        +
        Jmd
        +
        Jm^2
        +
        m^2d
        +
        m^3
    ).
}
\]

Дополнительная память:
\[
\boxed{
    O(md+m^2).
}
\]

\subsection{Рекомендуемая иерархия методов}

Для развития существующего truncated-SVD solver рекомендуется
сравнить следующие методы.

\paragraph{Method A: current greedy SVD.}

Последовательное добавление одной rank-one component с текущей
инициализацией из residual-gradient.

\paragraph{Method B: preconditioned greedy SVD.}

Тот же greedy solver, но следующая component инициализируется
из top singular pair
\[
    G^{-1/2}Q_t.
\]

\paragraph{Method C: isotropic truncated-SVD initialization + greedy refinement.}

Начальное приближение:
\[
    B_0
    =
    G^{-1/2}
    [G^{-1/2}C_B]_r.
\]

После этого используется существующий exact greedy refinement.

\paragraph{Method D: isotropic initialization + block ALS.}

Все $r$ компонент инициализируются одновременно через truncated SVD,
после чего совместно оптимизируются block ALS.

\paragraph{Method E: adaptive block spectral expansion.}

Rank начинает с малого $r_0$ и увеличивается блоками размера $q$
через
\[
    [G^{-1/2}Q_t]_q.
\]

После каждого expansion выполняется exact refinement.

\subsection{Главная проверяемая гипотеза}

Необходимо проверить, является ли preconditioned spectral matrix
\[
\boxed{
    G^{-1/2}C_B
}
\]
хорошей аппроксимацией структуры полного rank-constrained HPAO
решения.

Для полного unconstrained решения $B_\star$ и low-rank решения $B_r$
измерять:

\[
    F(B_r)-F(B_\star),
\]
\[
    \frac{\|B_r-B_\star\|_F}{\|B_\star\|_F},
\]
\[
    \text{projector distance после canonicalization},
\]
\[
    \text{число проходов через }U,
\]
\[
    \text{wall time}.
\]

Отдельно измерять точность isotropic model:
\[
\boxed{
    \theta
    =
    \frac{
        \sum_j
        N_j
        \|g_j\|_2^2
        \|H_j-\eta_jI\|_2
    }{
        \lambda
    }.
}
\]

Если
\[
    \theta\ll1,
\]
spectral initialization имеет строгую perturbation-интерпретацию.

Если
\[
    \theta\ge1,
\]
она остаётся допустимой эвристической инициализацией,
но multiplicative spectral certificate больше не применим.

\subsection{Итог}

Для pure truncated-SVD solver строго доказано следующее.

Для общего HPAO functional:

\[
\boxed{
    \text{лучший rank-one gain}
    =
    \|Q\|_{\mathrm{csp},\mathcal L}^2.
}
\]

Conditional directions удовлетворяют
\[
\boxed{
    a
    \propto
    G(v)^{-1}Qv,
    \qquad
    v
    \propto
    H(a)^{-1}Q^\top a.
}
\]

При точной изотропии:
\[
\boxed{
    B_r^\star
    =
    G^{-1/2}
    [G^{-1/2}C_B]_r.
}
\]

При этом
\[
\boxed{
    F(B_r^\star)-F(B_\star)
    =
    \sum_{k>r}s_k^2,
}
\]
а лучший следующий rank-one gain равен
\[
\boxed{
    s_{r+1}^2.
}
\]

Поэтому наиболее естественное развитие pure SVD solver имеет вид
\[
\boxed{
\text{whitened truncated-SVD initialization}
\rightarrow
\text{exact low-rank refinement}
\rightarrow
\text{preconditioned spectral residual expansion}.
}
\]

При этом сама оптимизируемая величина на всех этапах остаётся
\[
\boxed{B}
\]
и никакая correction-параметризация
$B=P+\Delta$ не используется.
